%% kapa: exam=1 round=35 year=2024 period=2 form=A subject=appraisal_law subject_ko=감정평가관계법규 seq=12 number=12
%% source: Q-Net 감정평가사 기출문제 (cst00309) artlSeq=5213229 file=raw/1008/5213229/2교시 문제지 원본.pdf
%% conversion: kordoc 4.24.51 -> tools/gen1.py (auto); review=proofread-sonnet; flags=none
\begin{question}{12}
국토의 계획 및 이용에 관한 법령상 시가화조정구역에 관한 설명으로 옳지 않은 것은?
\begin{choices}
  \item 시가화를 유보할 수 있는 기간은 5년 이상 20년 이내이다.
  \item 시가화조정구역의 지정에 관한 도시ㆍ군관리계획 결정이 있는 경우 결정 당시 이미 허가를 받아 공사에 착수한 자는 관할 관청에 신고하고 그 공사를 계속할 수 있다.
  \item 시가화조정구역에서 해제되는 구역 중 계획적인 개발 또는 관리가 필요한 지역에 대하여는 지구단위계획구역을 지정할 수 있다.
  \item 시가화조정구역에서 입목의 조림 또는 육림은 관할 관청에 신고하고 그 행위를 할 수 있다.
  \item 시가화조정구역의 지정에 관한 도시ㆍ군관리계획의 결정은 시가화 유보기간이 끝난 날의 다음날부터 그 효력을 잃는다.
\end{choices}
\end{question}
\answer{4}
